\documentclass{subfiles}
\begin{document}\label{Sec:6}
    \begin{definition*}
        Let \(\C* \subseteq A^{n}\) be a code. If \(\forall \bv{w} \in \C*\),
            the cyclic shift of \(\bv{w}\) is still a word of the code, 
            we say that \(\C*\) is cyclic.
    \end{definition*}

    \begin{remark*}
        Algebraically speaking, we say that \(\C* \subseteq A^{n}\) is cyclic
            if and only if \(\C*\) is an ideal of \(A^{n}\).
    \end{remark*}

    In our discussion of cyclic codes we treat words in a code as polynomials.
        That is, if \(\bv{w} = (w_{0}, w_{1}, \ldots, w_{k - 1})\) is a word of some cyclic code,
        then it refers to the polynomial \(f(x) = w_{0} + w_{1}x + \ldots + w_{k - 1}x^{k - 1}\).

    As for the case of linear codes, to each cyclic code we can associate two polynomials:
        one that generates the code, call it \(g(x)\), and one for code checking.
        It can be proved that \(g(x)\) is minimal and has degree \(n - k\),
        additionally it divedes the polynomial \(x^{n} - 1\).

    That is, a cyclic code \(\C*\) defines the set 
        \[
            \gen{g(x)} = \set*{f(x) \in \F_{n}[x]}[f(x) = g(x)u(x), u(x) \in \F_{n - \deg g}[x]].
        \]
    It should be noted that a cyclic shift in this context corresponds to a multiplication by \(x\).
        In fact, if \(\bv{w} = (w_{0}, w_{1}, \ldots, w_{k - 1})\) is a word, then 
        \((\bv{w}' = (w_{k - 1}, w_{0}, \ldots, w_{k - 2})\) is also a word.
        Hence, using the same matrix notation as for linear codes,
        the generator matrix for the code \(\C*\), with generator \(g(x)\), is
        \[
            G = \begin{pmatrix}
                g_{0}   & g_{1} & \ldots    & g_{n - k} &           &   & \\ 
                        & g_{0} & g_{1}     & \ldots    & g_{n - k} &   & \\
                        &       & \ddots     &\ddots     & \ldots   & \ddots & \\ 
                        &       &           & g_{0} & g_{1} & \ldots & g_{n - k} \\
            \end{pmatrix}
        \]

    On the other hand, to understand how the parity matrix for a cyclic code is obtained,
        we have to consider the following. Since \(\divides{g(x)}{x^{n} - 1}\), 
        and \(\deg g = n - k\), there must exist a polynomial \(h(x)\) such that 
        \[
            g(x)h(x) = x^{n} -1 \land \deg h = k.
        \]
        Though we don't descibe them, it should be noted that \(h(x)\) is the generating polynomial 
        for the dual code\footnotemark of \(\C*\). 

    Consider the following
        \[\begin{aligned}
            \gen{g(x)} \ni c(x) = f(x)g(x) & = c_{0} + c_{1}x + c_{2}x^{2} + \ldots + c_{n - 1}x^{n - 1} \\
            \gen{h(x)} \ni d(x) = l(x)h(x) & = d_{0} + d_{1}x + d_{2}x^{2} + \ldots + d_{n - 1}x^{n - 1} \\
        \end{aligned}\]
        if we multiply them we get 
        \[
            c(x)d(x) = f(x)g(x)l(x)h(x) = m(x)(x^{n} - 1) = m(x)x^{n} - m(x).
        \]
        Considering the degrees of \(f \text{ and } l\) we have:
        \[
            \deg f \le n - 1 - \deg g \le k - 1 \land \deg l \le n - 1 - \deg h \le n - k - 1
        \]
        that is, \(\deg m(x) \le n - 2\), meaning no term of degree \(n - 1\) exists.
        In other words, the sum 
        \[
            c_{0}d_{n - 1} + c_{1}d_{n - 2} + \cdots + c_{n - 1}d_{0} = 0;
        \]
        equivalently, the product \((c_{0}, \ldots, c_{n - 1}) * (d_{n - 1}, \ldots, d_{0})^{t} = 0\).

        If we consider it in is polynomial form, the second term of the product is nothing else but 
            the reverse of \(d(x)\), which is obliviously generated by some polynomial \(\tilde{h}(x)\)
            in turns reverse of \(h(x)\). 

        Summing up, the matrix \(H\) generated by \(\tilde{h}(x)\) is exactly the parity matrix for our code \(\C*\).

        \footnotetext{Informally speaking, these can be thought as the complement of some code \(\C*\).}
\end{document}
